\documentclass[twocolumn]{aastex631}

\usepackage{amsmath}
\usepackage{multirow}
\usepackage{natbib}
\usepackage{graphicx} 
\usepackage{aas_macros}

\begin{document}

Degenerate higher-order scalar-tensor theories are constructed to be free of the catastrophic Ostrogradsky ghost instability at the classical level by imposing specific algebraic constraints, but the quantum robustness of this mechanism is an open question. We address this by proposing a dynamical test: we treat the degeneracy conditions as defining a "healthy" manifold within the infinite-dimensional space of theories and investigate whether this manifold is stable under the Renormalization Group (RG) flow. Using the path integral formalism and the background field method, we compute the full one-loop beta-functionals for the general Type Ia class of these theories by calculating the UV divergences via heat kernel techniques. Our central result is that the RG flow vector is generically not tangent to the degeneracy manifold. This demonstrates that quantum fluctuations inevitably drive the theory away from the healthy subspace, dynamically resurrecting the ghost and its associated pathologies. We conclude that the classical stability mechanism is an unprotected fine-tuning, implying that this class of theories is fundamentally unstable at the quantum level without a yet-unknown protective symmetry.

\end{document}
                